%%% BEGIN fig:meson-production
\begin{figure}[htbp]
\centering
\begin{tikzpicture}[fd,
  t3 hadron/.style={draw=black!70,line width=1.5pt,line cap=butt},
  t3 lab/.style={font=\small},
  t3 note/.style={font=\footnotesize}]
  \node[hadron blob,ellipse,minimum width=3.0cm,minimum height=1.0cm,inner sep=0pt] (F) at (0,0) {GPD};
  \node[jetB blob,ellipse,minimum width=1.6cm,minimum height=0.8cm,inner sep=0pt] (M) at (0.9,4.3) {};
  \coordinate (Fl) at (-0.95,0.39);  \coordinate (Fr) at (0.95,0.39);
  \coordinate (G1) at (-0.95,1.7);   \coordinate (G2) at (0.95,1.7);
  \coordinate (V) at (-0.95,3.0);
  \coordinate (Mq) at (0.30,4.036);  \coordinate (Mb) at (1.20,3.93);
  \coordinate (Qs) at (-2.6,3.9);
  % hard kernel
  \draw[region] (-1.45,1.2) rectangle (1.42,3.45);
  \draw[t3 hadron] (-3.2,0) -- (F.west);  \draw[t3 hadron] (F.east) -- (3.2,0);
  \draw[t3 hadron] (M.east) -- (3.0,4.3);
  \draw[gluon] (G1) -- (G2);
  \draw[photon] (Qs) -- (V);
  \draw[quark] (Fl) -- (G1);  \draw[quark] (G1) -- (V);  \draw[quark] (V) -- (Mq);
  \draw[quark] (Mb) -- (G2);  \draw[quark] (G2) -- (Fr);
  \foreach \p in {G1,G2,V} \node[vertex] at (\p) {};
  \node[jetB blob,ellipse,minimum width=1.6cm,minimum height=0.8cm,inner sep=0pt] at (0.9,4.3) {};
  % labels
  \node[t3 lab,above left=-1pt] at (Qs) {$\gamma^*_L(q)$};
  \path[mom'={},mom sep=0.22cm,mom length=0.45cm] (Qs) -- (V);
  \path[mom={$(x+\xi)\bar P^+$},mom sep=0.26cm,mom length=0.45cm,mom pos=0.34] (Fl) -- (G1);
  \path[mom={$(x-\xi)\bar P^+$},mom sep=0.26cm,mom length=0.45cm,mom pos=0.66] (G2) -- (Fr);
  \path[mom={$zq'$},mom sep=0.24cm,mom length=0.42cm,mom pos=0.6] (V) -- (Mq);
  \path[mom'={$(1-z)q'$},mom sep=0.24cm,mom length=0.42cm,mom pos=0.86] (G2) -- (Mb);
  \path[mom'={$P$},mom sep=0.26cm,mom length=0.5cm] (-3.2,0) -- (F.west);
  \path[mom'={$P'$},mom sep=0.26cm,mom length=0.5cm] (F.east) -- (3.2,0);
  \node[t3 lab,right=1pt] at (3.0,4.3) {$M(q')$};
  \node[t3 lab,above=1pt] at (M.north) {$\phi_M(z)$};
  \node[t3 note,below=2pt] at (F.south) {$F^q(x,\xi,t)$};
  \node[t3 lab,anchor=east] at (-1.5,1.55) {$T_H$};
\end{tikzpicture}
\caption{硬遍举介子产生 $\gamma^*_L(q)+p(P)\to M(q')+p(P')$ 的一个领头阶图（夸克 GPD 的贡献；记号同\cref{fig:dvcs-handbag}）。
下方灰色团为 GPD 关联函数 $F^q(x,\xi,t)$，按介子的量子数取 $H^q,E^q$ 或 $\tilde H^q,\tilde E^q$；上方橙色团为介子的光锥分布振幅 $\phi_M(z)$，其中夸克带动量 $zq'$，反夸克带 $(1-z)q'$；
虚线框内为硬散射核 $T_H$，领头阶含一个硬胶子交换。振幅的因子化形式为
$\mathcal M=\sum_q\int_{-1}^{1}\dd x\int_0^1\dd z\,F^q(x,\xi,t)\,T^q_H(x,\xi,z;Q,\mu)\,\phi_M(z;\mu)+\text{幂次修正}$（Collins--Frankfurt--Strikman 1997），
只对纵向极化的 $\gamma^*$ 成立（全阶论证 [A]）。对横向极化光子，相应的扭度 3 卷积在端点 $z\to0,1$ 处发散，标准的共线因子化不成立。
对矢量介子，胶子 GPD 在同一阶也有贡献（图略）。}
\label{fig:meson-production}
\end{figure}
%%% END fig:meson-production
